\documentclass[9pt]{beamer}
\usepackage[utf8]{inputenc}
\usetheme{Madrid}
%%%%%%%%%%%%%%%%%%%%%%%%%%%%%%%%%%%%%%%%%%%%%%%%%%%%%%%%%%%%%%%%%%%%%%%%%%%%%%%%%%%%%%%%%%%%%%%%%%%%%%%%%%%%%%%%%%%%%%%%%%%%%%%%%%%%%%%
\definecolor{blue_text}{RGB}{0,115,197}
\definecolor{blue1}{RGB}{26,129,203}
\definecolor{blue2}{RGB}{102,171,220}
\definecolor{blue3}{RGB}{154,199,232}

\definecolor{green1}{RGB}{0,128,0}
\definecolor{green2}{RGB}{204,229,204}
\definecolor{green3}{RGB}{229,242,229}

\definecolor{lightblue2}{RGB}{204,227,243}
\definecolor{lightblue3}{RGB}{229,241,249}

\definecolor{orange0}{RGB}{255,116,0}
\definecolor{blue0}{RGB}{0,103,169}%vaidya
\definecolor{green0}{RGB}{0,148,104}


\setbeamercolor*{palette tertiary}{bg=blue1,fg=white}
\setbeamercolor*{palette secondary}{bg=blue2,fg=white}
\setbeamercolor*{palette primary}{bg=blue3,fg=black}

\setbeamercolor{titlelike}{bg=white,fg=blue_text}
\setbeamertemplate{title page}[default][rounded=false, shadow=false]

\setbeamercolor{frametitle}{bg=white,fg=blue_text}

\setbeamercolor{item projected}{bg=blue_text}

\setbeamertemplate{blocks}[rounded]%[shadow]
\setbeamercolor{block title example}{bg=green2,fg=green1}
\setbeamercolor{block body example}{bg=green3,fg=black}

\setbeamercolor{block title}{bg=lightblue2,fg=blue_text}
\setbeamercolor{block body}{bg=lightblue3,fg=black}

\setbeamercolor{button}{bg=blue2,fg=white}

\setbeamertemplate{navigation symbols}{}
%%%%%%%%%%%%%%%%%%%%%%%%%%%%%%%%%%%%%%%%%%%%%%%%%%%%%%%%%%%%%%%%%%%%%%%%%%%%%%%%%%%%%%%%%%%%%%%%%%%%%%%%%%%%%%%%%%%%%%%%%%%%%%%%%%%%%%%

\usepackage{eqnarray,amsmath}
\usepackage{amsfonts}
\usepackage{amssymb}
\usepackage{graphicx}
\usepackage{lmodern} 
\usepackage{bm}
\usepackage{epstopdf}
\usepackage{changepage}
\usepackage{array,booktabs}

\usepackage{pgfplots}
\usepackage{tikz}
\usetikzlibrary{patterns}
\usetikzlibrary{plotmarks}
\usepgfplotslibrary{fillbetween}

\DeclareMathOperator*{\argmax}{arg\,max}
\newcommand*\diff{\mathop{}\!\mathrm{d}}
\DeclareMathOperator{\supp}{supp}

\usepackage{transparent}
\newcommand{\semitransp}[2][35]{\color{fg!#1}#2}

\usetikzlibrary{arrows.meta}

%%%%%%%%%%%%%%%%%%%%%%%%%%%%%%%%%%%%%%%%%%%%%%%%%%%%%%%%%%%%%%%%%%%%%%%%%%%%%%%%%%%%%%%%%%%%%%%%%%%%%%%%%%%%%%%%%%%%%%%%%%%%%%%%%%%%%%%


%====================================================
%========== DEFINITION OF AUTHORS ETC...=============
%====================================================

\author[Ash, Lou, Song]{%
    \texorpdfstring{%
    \begin{columns}
        \column{.3333\linewidth}
        \centering
        Elliot Ash \\\vspace{.5em} \small ETH Zurich
        \column{.3333\linewidth}
        \centering
        Yichuan Lou \\\vspace{.5em} \small UTokyo
        \column{.3333\linewidth}
        \centering
        Lena Song \\\vspace{.5em} \small UIUC
    \end{columns}}{Elliot Ash, Yichuan Lou, Lena Song}
}

\title[AI and Communication]{AI and Communication}
%\subtitle{Subtitle}
\date[Conference Name]{May 2024}




%====================================================
%========== BEGINNING OF DOCUMENT ===================
%====================================================
\begin{document}

\begin{frame}[plain]
    \titlepage
\end{frame}

\begin{frame}{Compositions of Senders and Recipients}

\begin{itemize}\setlength{\itemsep}{1em}
    \item Both senders and recipients consist of three \textcolor{orange0}{identity} groups: humans, Gen-AI bots, and rule-based bots
    \item Denote by $m_h,m_g,m_r$ the number of different groups of senders 
    \begin{itemize}
        \item Let $\vec{m} \triangleq (m_h,m_g,m_r)$
    \end{itemize}
    \item Denote by $n_h,n_g,n_r$ the number of different groups of recipients
    \begin{itemize}
        \item Let $\vec{n} \triangleq (n_h,n_g,n_r)$
    \end{itemize}
\end{itemize}
\end{frame}

    

\begin{frame}{Human Sender}
\begin{itemize}\setlength{\itemsep}{1em}
    \item We study the behavior of \textcolor{orange0}{human senders} who choose the level of their contribution in some prosocial activity \textcolor{red}{(YL: can be more specific later such as posting content)}
    \item We assume a human sender has two types, high $(\theta^+)$ and low $(\theta^-)$
    \item Each human sender chooses a contribution level $a\in \mathbb{R}$, at cost $C(a)$
    \begin{itemize}
        \item Denote by $a^+$ ($a^-$) the action chosen by type $\theta^+$ ($\theta^-$)
    \end{itemize}
    \item A human sender's utility depends on the three key ingredients: intrinsic motivation, social interaction incentives, and reputational concerns:
    \begin{equation*}
        U = (\theta + \lambda )a - C(a) + \gamma \mathbb{P}(human|a)\mathbb{P}(\theta=\theta^+|a),
    \end{equation*}
    \begin{itemize}
        \item Let $\mu\triangleq \gamma \mathbb{P}(human|a)$, where $\mathbb{P}(human|a)$ is a recipient's updated belief that the sender is a human
        \item We will show $\mathbb{P}(human|a)$ can be set independent of $a$; so $\mu$ can be treated as an exogenous parameter
    \end{itemize}
\end{itemize}
\end{frame}

\begin{frame}{Human Sender: Intrinsic Utility}
\begin{equation*}
    U = (\theta + \lambda )a - C(a) + \mu \mathbb{P}(\theta=\theta^+|a),
\end{equation*}
\vspace{.5em}
\begin{itemize}\setlength{\itemsep}{1em}
    \item The term $\theta a$ captures the intrinsic utility from contribution
    \item Intrinsic utility assumes senders derive direct utility from contribution, and leads to ``the doing of an activity for its inherent satisfactions rather than for some separable consequence"
\end{itemize}
\end{frame}


\begin{frame}{Human Sender: Social Interaction Utility}
\begin{equation*}
    U = (\theta + \lambda )a - C(a) + \mu \mathbb{P}(\theta=\theta^+|a),
\end{equation*}
\vspace{.5em}

\begin{itemize}\setlength{\itemsep}{1em}
    \item The term $\lambda a$ captures the social interaction utility from contribution
    \item The parameter $\lambda$ measures the intensity of social interaction incentives
    \item The social interaction utility assumes senders derive utility from interacting with different groups of recipients 
    \item We assume $\lambda$ is increasing in $n_h,n_g,m_h,m_g$
    \vspace{.5em}
    \begin{itemize}
        \item E.g., $\lambda = (n_h+n_g)(m_h+m_g)$
    \end{itemize}
\end{itemize}
\end{frame}

\begin{frame}{Human Sender: Image-Related Utility}
\begin{equation*}
    U = (\theta + \lambda )a - C(a) + \mu \mathbb{P}(\theta=\theta^+|a),
\end{equation*}
\vspace{.5em}

\begin{itemize}\setlength{\itemsep}{1em}
    \item The term $\mu \mathbb{P}(\theta=\theta^+|a)$ captures the reputational utility from contribution
    \item The parameter $\mu$ measures the intensity of image concerns
    \item We assume $\mu$ is increasing in $n_h$
    \vspace{.5em}
    \begin{itemize}
        \item E.g., $\mu = n_h \mathbb{P}(human)$
    \end{itemize}
\end{itemize}
\end{frame}

\begin{frame}{Bot Sender}
\begin{itemize}\setlength{\itemsep}{1em}
    \item A bot sender, either Gen-AI or rule-based, has no types or preferences
    \item We assume each Gen-AI or rule-based bot sender chooses some fixed actions:
\end{itemize}
\vspace{1em}

\begin{exampleblock}{Assumption 1}
    A rule-based bot sender always takes a sufficiently small action $a_r$.
\end{exampleblock}

\begin{exampleblock}{Assumption 2}
    A Gen-AI sender always takes a fixed action $a_g$. We consider two different cases:
    \begin{itemize}
        \item Case 1: A Gen-AI bot sender always takes a non-equilibrium action, i.e., $a_g\neq a^+,a^-$.
        \item Case 2: A Gen-AI bot sender always mimic a high type human sender's equilibrium action, i.e., $a_g=a^+$.
    \end{itemize}
\end{exampleblock}
\end{frame}


\begin{frame}{Signaling Game}

Timing:
\begin{itemize}\setlength{\itemsep}{1em}
    \item A single sender is matched with recipients $(m_h,m_g,m_r)$ 
    \item The sender's identity (human or bot) and type ($\theta\in \{\theta^+,\theta^-\}$ if human) is initially unknown to the recipients
    \item The sender chooses an action $a\in \mathbb{R}$
    \item The belief about the sender's identity and type is updated
    \item If the sender is a human, his payoff is realized
\end{itemize}

\end{frame}



\begin{frame}{Unique Equilibrium}\label{prop1}
\begin{exampleblock}{Proposition 1}
    Fix any positive parameters $\lambda$ and $\mu$. The unique PBE surviving the D1 criterion (Cho and Kreps, 1987) is the most efficient separating equilibrium, in which there exists a threshold $\theta^*\triangleq\tau(\theta^-,\mu,\lambda)$ such that
    \begin{align*}
    \begin{cases}
        a^+ = C'^{-1}(\theta^++\lambda), &a^- = C'^{-1}(\theta^-+\lambda)  \quad\text{if } \theta^+\geq \theta^*\\ 
        a^+ = C'^{-1}(\theta^*+\lambda), &a^- = C'^{-1}(\theta^-+\lambda)     \quad\text{otherwise.}
    \end{cases}
    \end{align*}
    Moreover, the threshold function $\tau(\cdot)$ is increasing in both $\theta^-$ and $\mu$. If $C'''(\cdot)>0$, then $\tau(\cdot)$ is increasing in $\lambda$; if $C'''(\cdot)<0$, then $\tau(\cdot)$ is decreasing in $\lambda$.
\end{exampleblock}
\vspace{1em}

Remark:
\begin{itemize}\setlength{\itemsep}{.5em}
    \item When $\theta^+\geq \theta^*$, both equilibrium actions $a^+$ and $a^-$ are increasing in social interaction incentives $(\lambda)$ but independent of reputational incentives ($\mu$)
    \item When $\theta^+<\theta^*$, $a^+$ is increasing in $\mu$ for sure, whereas $a^-$ is increasing in $\lambda$ but independent of $\mu$
    \item \textcolor{red}{(YL: how equilibrium outcomes rely on $\lambda$ and $\mu$ will allow us to study the effect of AI)}
\end{itemize}

\vspace{.5em}
\hyperlink{proof_prop1}{\beamergotobutton{Proof}}

\end{frame}


\begin{frame}{Example: $C(a) = \frac{1}{2}a^2$}

\begin{itemize}\setlength{\itemsep}{1em}
    \item In this case, $\tau(\theta^-,\mu,\lambda) = \theta^- + \sqrt{2\mu}$
    \item The equilibrium contribution levels are given by
    \begin{align*}
    \begin{cases}
        a^+ = \theta^++\lambda,  &a^- = \theta^-+\lambda  \quad\text{if } \theta^+\geq \theta^-+\sqrt{2\mu} \\
        a^+ = \theta^-+\lambda+\sqrt{2\mu}, &a^- =  \theta^-+\lambda  \quad \text{otherwise.}
    \end{cases}
    \end{align*}
\end{itemize}
\end{frame}

\begin{frame}{What's Next?}
\begin{itemize}\setlength{\itemsep}{1em}
    \item So far we have treated $\mu$ and $\lambda$ as given and have learned how they affect equilibrium outcomes
    \item In the following, we will conjecture and analyze how the composition of senders $\vec{m} = (m_h,m_g,m_r)$ and the composition of recipients $\vec{n} = (n_h,n_g,n_r)$ determine the values of $\mu$ and $\lambda$, and thus indirectly affect equilibrium outcomes
\end{itemize}
\end{frame}

\begin{frame}{What determines $\mu$?}
\begin{itemize}\setlength{\itemsep}{1em}
    \item Recall that Assumption 2 considers two cases:
    \begin{itemize}
        \item In Case 1, a Gen-AI sender chooses a fixed non-equilibrium action $a_g\neq a^+$
        \item In Case 2, a Gen-AI sender mimics a high type human sender, i.e., $a_g= a^+$
    \end{itemize}
    \item In Case 1, upon observing $a^+$, a recipient knows the sender must be human, hence $\mathbb{P}(human|a^+)=1$, which implies 
    \begin{align*}
        \mu_1 \triangleq \gamma \mathbb{P}(human|a^+) = \gamma
    \end{align*}
    \item In Case 2, upon observing $a^+$, a recipient only knows the sender cannot be a rule-based bot, hence $\mathbb{P}(human|a^+)=\frac{m_h}{m_h+m_g}$, which implies 
    \begin{align*}
        \mu_2\triangleq \gamma \mathbb{P}(human|a^+) = \gamma \frac{m_h}{m_h+m_g}
    \end{align*}
    \item An immediate implication is the presence of Gen-AI senders may reduce a human sender's incentive to invest as it may reduce $a^+$ 
\end{itemize}
\end{frame}


\begin{frame}{What determines $\mu$?}
As for the parameter $\gamma$, we assume it only depends on the number of human recipients. That is,
\begin{align*}
    \gamma \triangleq n_h.
\end{align*}
\end{frame}

\begin{frame}{What determines $\lambda$?}
Most general form: $\lambda \triangleq \lambda(\vec{n},\vec{m})$
\end{frame}




\begin{frame}{Proof of Proposition 1}\label{proof_prop1}

Let $p\triangleq \mathbb{P}(\theta=t_h)$. The sender utility function becomes
\begin{align*}
    U(\theta,a,p)  \triangleq (\theta+\lambda)a-C(a)+p
\end{align*}
Define $U^*(\theta^-) \triangleq U(\theta^-,a^-,0)$ and $U^*(\theta^+) \triangleq U(\theta^+,a^+,1)$ such that $a^-$ solves 
\begin{align*}
    \max_{a} U(\theta^-,a,0)
\end{align*}
and $a^+$ solves
\begin{align*}
    \max_{a} U(\theta^+,a,1)\quad \text{subject to} \quad U(\theta^-,a^+,1)\leq U^*(\theta^-)
\end{align*}
It is easy to see that the pair $(a^+,a^-)$ corresponds exactly to the form specified in Proposition 1.
\vspace{1em}

The above process produces a signaling strategy for the sender and a belief rule for the recipient. To complete the belief rule, assume that the recipient's belief $p = 0$ for $a<a^+$ and $p=1$ for $a\geq a^+$. It is easy to check the solution is the most efficient separating equilibrium, or the ``Riley outcome". In our setup, it is the unique equilibrium satisfying D1 criterion (Cho and Kreps, 1987). 
\end{frame}


\begin{frame}{Proof of Proposition 1 (Cont'd)}

The threshold function $\tau(\theta^-,\mu,\lambda)$ is defined by the indifference condition $U^*(\theta^-)= U(\theta^-,\hat{a},1)$ where $\hat{a}$ solves $\max_aU(\tau(\theta^-,\mu,\lambda),a,1)$, or equivalently, $C'(\hat{a}) = \tau(\theta^-,\mu,\lambda)+\lambda$. The indifference condition becomes
\begin{align*}
    (\theta^-+\lambda)a^- - C(a^-) = (\theta^-+\lambda)\hat{a} - C(\hat{a})+\mu 
\end{align*}
By the implicit function theorem, 
\begin{align*}
    \frac{\partial \tau(\theta^-,\mu,\lambda)}{\partial \theta^-} &= \frac{\hat{a}- a^-}{(\tau(\theta^-,\mu,\lambda) - \theta^-)\hat{a}'}>0,\\
    \frac{\partial \tau(\theta^-,\mu,\lambda)}{\partial \mu} &= \frac{1}{(\tau(\theta^-,\mu,\lambda) - \theta^-)\hat{a}'}>0,\\
    \frac{\partial \tau(\theta^-,\mu,\lambda)}{\partial \lambda} &= -\frac{a^--\hat{a} +(\tau(\theta^-,\mu,\lambda)-\theta^-)\hat{a}'}{(\tau(\theta^-,\mu,\lambda) - \theta^-)\hat{a}'}.
\end{align*}
By Taylor's theorem, we have 
\begin{align*} 
    \begin{cases}
        \frac{\partial \tau(\theta^-,\mu,\lambda)}{\partial \lambda}>0\quad &\text{if }C'''(\cdot)>0,\\
        \frac{\partial \tau(\theta^-,\mu,\lambda)}{\partial \lambda}<0\quad &\text{if }C'''(\cdot)<0,\\
        \frac{\partial \tau(\theta^-,\mu,\lambda)}{\partial \lambda}=0\quad &\text{if }C'''(\cdot)=0.
    \end{cases}
\end{align*}

\vspace{1em}
\hyperlink{prop1}{\beamerreturnbutton{Back}}
	
\end{frame}














\end{document}
